\chapter{State of the Art}
\label{chap:sota}

\draftnote{Bibliography expanded to 64 verified references. Every entry cited in this chapter was checked against a primary record (ACL Anthology, PMLR, NeurIPS, ACM, JSTOR, EUR-Lex, or the publisher of record). Numerical results attributed to prior work are quoted from the cited papers and have not been reproduced here.}

\section{Scope and Organisation}
\label{sec:sota-scope}

This chapter reviews the research that the present work builds upon and, more importantly, the research that defines what would count as an adequate answer to the problem it addresses. The review is organised around the pipeline under study rather than around chronology. Section~\ref{sec:sota-rag} establishes why an external evidence store is required at all, and Section~\ref{sec:sota-parametric} examines the competing option of adapting model weights instead. Sections~\ref{sec:sota-retrieval} to~\ref{sec:sota-rerank} treat the retrieval stack: lexical, dense, and hybrid candidate generation, the choice of indexing unit, and cross-encoder reranking. Section~\ref{sec:sota-legal} turns to the specific properties of legal and regulatory language and to the empirical record of automated systems operating on it. Sections~\ref{sec:sota-attribution} and~\ref{sec:sota-selective} cover attribution and selective answering, which together constitute the safety layer of the architecture. Sections~\ref{sec:sota-eval} and~\ref{sec:sota-validity} address measurement: what to measure, and how to know that the measuring instrument is itself sound. Section~\ref{sec:sota-agentic} treats multi-step and agentic retrieval, and Section~\ref{sec:sota-gap} states the gap that this work occupies.

Two commitments shape the selection of material. First, a claim about system performance is only admitted here if the cited work reports how it was measured; several widely repeated assertions about retrieval-augmented generation do not survive that filter. Second, work that documents \emph{failure} is treated as at least as important as work that reports improvement, because the contribution of this project is a policy for responding to failure rather than a new architecture.

\section{From Parametric Knowledge to Retrieval-Augmented Generation}
\label{sec:sota-rag}

Transformer architectures made it possible to model long-range dependencies through attention rather than recurrence~\cite{vaswani2017attention}, and bidirectional pre-training subsequently improved language understanding and transfer across many \ac{NLP} tasks~\cite{devlin2019bert}. Knowledge acquired this way is stored implicitly in model weights. That storage mechanism has three properties that are unacceptable in a regulatory setting: it is not addressable, so a statement cannot be traced to the text that licensed it; it is not updatable, so an amended provision cannot be corrected without retraining; and it is not auditable, so a reader cannot verify a claim without independently rediscovering its source.

\ac{RAG} addresses these properties by conditioning generation on passages retrieved from an external collection at inference time. Lewis et al.\ jointly modelled retrieval and sequence generation for knowledge-intensive tasks~\cite{lewis2020rag}, and REALM integrated retrieval directly into language-model pre-training so that the retriever was learned rather than fixed~\cite{guu2020realm}. Subsequent architectures explored how retrieved evidence should enter the generator. Fusion-in-Decoder encodes each passage independently and allows the decoder to attend jointly across all of them, with accuracy improving as more passages are supplied~\cite{izacard2021fid}. RETRO takes the scaling argument further, showing that a 7.5-billion-parameter model with chunked cross-attention over a two-trillion-token retrieval database can match models an order of magnitude larger~\cite{borgeaud2022retro}, which establishes that capability can be relocated from weights into an index. REPLUG inverts the usual adaptation target by treating the language model as a frozen black box and training the retriever against its feedback~\cite{shi2024replug}; this matters here because the generator in the present system is a commercial \ac{API} that cannot be fine-tuned, so all available adaptation is on the retrieval side.

The literature has since consolidated into a recognised taxonomy of naive, advanced, and modular \ac{RAG}, decomposing the pipeline into indexing, retrieval, augmentation, and generation stages each with characteristic failure modes~\cite{gao2023ragsurvey}. That decomposition is adopted throughout this work. It is worth stating plainly that ``\ac{RAG}'' is not a system but a family of separable decisions: corpus construction, chunking unit, retrieval model, candidate depth, reranking, context assembly, prompt contract, and verification policy. Reporting a single accuracy figure for ``\ac{RAG}'' without specifying these choices communicates very little, and comparisons across papers that differ on several of them simultaneously are not informative about any one of them.

Adaptive retrieval marks the current frontier. Self-RAG trains a model to emit reflection tokens deciding when retrieval is necessary and then to critique whether its own output is supported by the retrieved evidence, improving both factuality and citation precision over fixed-retrieval baselines~\cite{asai2024selfrag}. This is the closest published antecedent to the escalation policy developed in Chapter~\ref{chap:method}, and the difference is instructive: Self-RAG learns the control decision inside the model, whereas the present work externalises it into an inspectable policy over observable retrieval signals, precisely so that the decision can be audited by a compliance reviewer.

In a high-stakes regulatory setting the objective is stricter than general factual usefulness. The selected evidence must identify the governing legal unit; each material claim must be entailed by that evidence; and the system must recognise when a question is ambiguous or when no retrieved source supports a safe answer.

\section{Retrieval or Parametric Adaptation?}
\label{sec:sota-parametric}

The alternative to retrieval is to write domain knowledge into the weights. Parameter-efficient methods have made this affordable. Low-Rank Adaptation freezes the pre-trained weights and trains only injected low-rank matrices, reducing trainable parameters by several orders of magnitude with no additional inference latency~\cite{hu2022lora}, and quantised adaptation backpropagates through a frozen four-bit model into those adapters, bringing the fine-tuning of very large models within reach of a single consumer \ac{GPU}~\cite{dettmers2023qlora}. The baseline experiment reported in Chapter~\ref{chap:results} uses exactly this route, and it is therefore not open to this work to dismiss fine-tuning on cost grounds.

The comparison must instead be made on knowledge properties. Ovadia et al.\ compare unsupervised fine-tuning against retrieval for injecting new and time-shifted knowledge across knowledge-intensive benchmarks and find retrieval consistently superior, with fine-tuning struggling in particular on facts absent from pre-training~\cite{ovadia2024finetuning}. For a regulatory corpus the argument is sharper still. Regulation~(EU)~No~376/2014 has been amended since adoption~\cite{eu2014reg376}, and the consolidated presentation used here is a dated revision~\cite{easa2022easyaccess376}; an adapter trained on one revision silently encodes it. More decisively, adapter weights cannot supply provenance. A fine-tuned model may produce a correct obligation and remain unable to identify which paragraph imposes it, which is the property the application actually requires.

The correct reading of the three-system baseline in this thesis is therefore not that fine-tuning is weak, but that the comparison is between architectures with different \emph{knowledge-access} properties, only one of which supports citation. This framing is stated here because the baseline results are otherwise open to a misreading that the earlier project documentation did not adequately guard against.

\section{Sparse, Dense, and Hybrid Retrieval}
\label{sec:sota-retrieval}

Sparse retrieval represents queries and documents through lexical statistics. BM25 remains a strong baseline because it rewards exact term overlap while normalising for document length~\cite{robertson2009bm25}. This is valuable in regulations, where article numbers, acronyms, defined terms, and formulaic expressions carry substantial meaning: a query naming Article~4(6) is not well served by a model that treats the numeral as noise. Learned sparse retrieval shows that the lexical channel is itself improvable; SPLADE learns sparse, expansion-aware term weights in the vocabulary space of a pre-trained encoder, retaining inverted-index efficiency while mitigating vocabulary mismatch~\cite{formal2021splade}. The vocabulary-mismatch weakness that motivates a dense channel is therefore a property of BM25 specifically rather than of sparse retrieval as such, and this qualification is carried through the experimental discussion.

Dense retrieval maps queries and passages into a learned vector space. Dense Passage Retrieval demonstrated that independently encoded questions and passages support efficient open-domain retrieval~\cite{karpukhin2020dpr}, and Sentence-BERT showed how siamese encoders make large-scale semantic similarity search computationally practical~\cite{reimers2019sentencebert}. Contriever established that contrastive pre-training on unlabelled text alone yields a dense retriever competitive with BM25 across the BEIR suite, and that it transfers across languages~\cite{izacard2022contriever}; this supports the use of a general-purpose commercial embedding model as a zero-shot first-stage retriever without in-domain training. Between these single-vector models and full cross-attention lies late interaction: ColBERT encodes query and document tokens separately and defers their comparison to a cheap maximum-similarity operator~\cite{khattab2020colbert}, and ColBERTv2 adds distillation and residual compression to cut the index footprint substantially while improving quality~\cite{santhanam2022colbertv2}. Late interaction is not used in this work; it is recorded here as a deliberate omission, traded away for a simpler index and no bespoke training.

Multilingual embedding models bear directly on the scope question raised in Chapter~\ref{chap:context}. M3-Embedding produces dense, lexical, and multi-vector representations across more than one hundred languages from a single model, and reports that combining its dense and sparse scoring modes outperforms either alone~\cite{chen2024m3embedding}. That internal result is independent corroboration of the premise underlying hybrid retrieval. Model selection across the embedding market is now conventionally argued through the Massive Text Embedding Benchmark, which showed that no single embedding model dominates across task families~\cite{muennighoff2023mteb} --- a finding that justifies selecting an embedding model against the target task rather than against a general leaderboard position.

Hybrid retrieval treats lexical and semantic evidence as complementary. Reciprocal-rank fusion combines ranked lists without requiring their raw scores to be calibrated~\cite{cormack2009rrf}. For a document $d$ retrieved by systems $i$, the fused score is

\begin{equation}
  \operatorname{RRF}(d) = \sum_i \frac{1}{c + \operatorname{rank}_i(d)},
  \label{eq:rrf}
\end{equation}

where $c$ is a fixed rank constant. The method is attractive for controlled experiments because it has a small and interpretable parameter surface and requires no score normalisation.

That attractiveness must be qualified. Bruch et al.\ analyse fusion functions for hybrid retrieval theoretically and empirically, and report that a convex combination of normalised scores outperforms \ac{RRF} both in and out of domain, noting that a rank-based method necessarily discards score magnitude and is sensitive to its rank constant~\cite{bruch2024fusion}. The present work nonetheless adopts equal-weight \ac{RRF} with a fixed constant, and the justification is methodological rather than performance-based: a fusion function with no tuned parameters cannot leak information from the held-out split, which matters more here than the last increment of recall. This trade-off is stated explicitly in Chapter~\ref{chap:method} and revisited among the limitations.

\begin{table}[htbp]
  \centering
  \caption{Retrieval families and their principal trade-offs, as characterised by the cited work.}
  \label{tab:retrieval-families}
  \small
  \begin{tabularx}{\textwidth}{@{}l l X@{}}
    \toprule
    \textbf{Family} & \textbf{Representative} & \textbf{Property relevant to regulatory text} \\
    \midrule
    Lexical sparse & BM25~\cite{robertson2009bm25} & Rewards exact identifiers and formulaic wording; blind to paraphrase. \\
    Learned sparse & SPLADE~\cite{formal2021splade} & Adds term expansion while remaining servable on an inverted index. \\
    Single-vector dense & DPR~\cite{karpukhin2020dpr}, Contriever~\cite{izacard2022contriever} & Recovers paraphrase; may underweight exact article references. \\
    Late interaction & ColBERT~\cite{khattab2020colbert,santhanam2022colbertv2} & Higher accuracy at greater index cost; not adopted here. \\
    Multilingual multi-mode & M3-Embedding~\cite{chen2024m3embedding} & Single model spanning languages and scoring modes. \\
    Rank fusion & \ac{RRF}~\cite{cormack2009rrf} & No score calibration, no tuned parameters; discards magnitude~\cite{bruch2024fusion}. \\
    \bottomrule
  \end{tabularx}
\end{table}

\section{Evidence Granularity and Document Hierarchy}
\label{sec:sota-granularity}

The choice of indexing unit is an independent lever on retrieval quality rather than an implementation detail. Chen et al.\ compare passage-, sentence-, and proposition-level indexing and find the finest granularity consistently strongest, with the largest gains on entities that are rare in the training distribution~\cite{chen2024densex}. Regulatory text amplifies this effect, because the normative content of a provision is typically carried by a single subparagraph while the surrounding article supplies scope and definitions. A coarse chunk containing several obligations dilutes the signal for each of them.

Fine granularity introduces the opposite hazard. A subparagraph detached from its heading, its defining article, and its table header may be uninterpretable, and in a legal document it may be actively misleading --- an exception read without its governing rule inverts the obligation. A parent--child design addresses this tension by retrieving atomic children while retaining links to coherent legal parents, so that context can be restored after ranking rather than before it.

Context restoration must nevertheless be bounded, and the reason is empirical. Liu et al.\ document a U-shaped performance curve in long contexts: models use evidence reliably at the beginning and end of the supplied window but degrade sharply when the relevant passage sits in the middle, with performance falling as the context grows even for models advertised as long-context~\cite{liu2024lostmiddle}. Two consequences follow for this work. First, the alternative design of supplying the entire regulation to a long-context model is not a safe default. Second, ordering and context length are substantive variables, which is what makes reranking worth its cost and what motivates a tight evidence budget rather than a generous one.

This project therefore treats hierarchy as explicit metadata rather than as a licence to promote every child of a long parent. Parent-aware expansion is applied only after retrieval and within a bounded context budget. The distinction matters because hierarchy can improve interpretability without necessarily improving relevance ranking, and the two effects must be measured separately.

\section{Cross-Encoder Reranking}
\label{sec:sota-rerank}

Bi-encoder retrieval is efficient because document vectors can be pre-computed and cached, but query and passage interact only through a similarity function applied to representations formed in ignorance of one another. A cross-encoder instead processes the pair jointly, allowing full attention between query and candidate terms. BERT-based passage reranking demonstrated substantial gains from this direct interaction~\cite{nogueira2019reranking}, and monoT5 showed that the effect is not an artefact of encoder-only models by recasting ranking as sequence-to-sequence target-token generation, with the additional finding that sequence-to-sequence rerankers degrade more gracefully in low-data regimes~\cite{nogueira2020monot5}. Reranking is therefore an architectural stage rather than a property of one model family.

The cost of full attention is tolerable only over a small candidate pool, which is what makes distillation relevant. MiniLM compresses a large transformer by matching the teacher's self-attention distributions and value relations, producing students that retain the large majority of teacher performance at a fraction of the cost~\cite{wang2020minilm}. The reranker used in this work is a MiniLM cross-encoder trained on MS MARCO, a large collection of anonymised web search queries with human-generated answers and passage judgements~\cite{nguyen2016msmarco}. Naming that training distribution is not a formality. MS MARCO queries are short, factoid-oriented, and drawn from general web search; questions about regulatory obligations are longer, structurally explicit, and drawn from a specialist sublanguage. Any gain observed from this reranker is therefore obtained \emph{despite} a domain shift, and any residual failure must be interpreted with that shift in view.

Reranking and candidate generation solve different problems, and conflating them is a recurrent source of confusion in applied reports. A reranker can move a relevant passage from rank~18 to rank~3; it cannot recover a passage absent from the candidate pool. Candidate recall at the fusion depth is thus a hard ceiling on the reranked result, and the two quantities must be reported separately. This distinction is central to the switching policy developed in this work.

\section{Legal and Regulatory Language Processing}
\label{sec:sota-legal}

Legal text is a genuine distribution shift rather than ordinary domain vocabulary. LEGAL-BERT showed that domain-adaptive pre-training on legislation and case law yields consistent downstream gains, and that the standard fine-tuning recipe does not transfer unmodified to specialist domains~\cite{chalkidis2020legalbert}. CaseHOLD refined the condition under which such adaptation pays: gains materialise when the task is both sufficiently domain-specific and sufficiently difficult, and are marginal otherwise~\cite{zheng2021casehold}. LexGLUE subsequently standardised evaluation across seven English legal datasets, including tasks defined over EU legislation~\cite{chalkidis2022lexglue}.

These benchmarks are informative about classification and, read carefully, they are also informative about what is missing. They score aggregate accuracy over labels; none of them asks whether a generated statement is attributable to a specific provision. LegalBench makes the gap visible from a different direction. Built collaboratively by legal professionals around a typology of legal reasoning, it shows that model performance is highly uneven across reasoning types, with issue-spotting comparatively strong and rule application and conclusion comparatively weak~\cite{guha2023legalbench}. A single accuracy figure therefore conceals a task-dependent competence profile, which is precisely the information a compliance deployment requires in order to decide where a human must remain in the loop. SARA makes the same point on statutory material specifically: pairing tax statutes with entailment and numerical cases, it shows that models trained on general corpora apply rule text to fact patterns poorly~\cite{holzenberger2020sara}. Questions over aviation regulation are rule-application questions of exactly this kind, so retrieval quality cannot be assumed to entail correct application.

The most directly relevant empirical result concerns deployed systems. Magesh et al.\ audit commercial retrieval-augmented legal research tools and find that grounding reduces but does not eliminate hallucination, with error rates in the range of roughly one in six to one in three responses, including miscited and unsupported authority~\cite{magesh2025hallucination}. Their distinction between an answer that is wrong and an answer that is right but misgrounded is adopted directly in this thesis. It is the empirical basis for the claim that retrieval grounding is not by itself a safety property, and that a citation-bearing system must be evaluated on groundedness rather than on a headline accuracy figure.

Aviation regulatory text has attracted its own resources. aeroBERT constructed an annotated corpus of aerospace requirements drawn from certification regulations and showed that a domain-adapted encoder substantially outperforms general-purpose zero-shot classifiers on it~\cite{tikayatray2023aerobert}, establishing the precedent for treating aviation regulation as a distinct \ac{NLP} domain with its own resources. Related work on occurrence narratives shows that free-text safety reports carry signal absent from structured fields~\cite{miyamoto2022aviation}. That literature addresses the reports produced \emph{under} an occurrence-reporting regime; the present work addresses the regime itself, that is, the regulation that determines what must be reported, by whom, and within what deadline~\cite{eu2014reg376,easa2022easyaccess376}. To the author's knowledge no published benchmark targets citation-bearing question answering over this instrument, which is the concrete gap the corpus and benchmark of Chapter~\ref{chap:method} are built to fill.

\section{Attribution, Citation, and Verifiability}
\label{sec:sota-attribution}

Hallucination is standardly divided into intrinsic cases, where output contradicts the source, and extrinsic cases, where output is simply unverifiable against it~\cite{ji2023hallucination}; more recent surveys reframe the same territory as a distinction between factuality and faithfulness and trace causes through data, training, and inference~\cite{huang2025hallucination}. In a regulatory application the extrinsic case dominates the risk profile, because a fluent statement that no provision supports is both plausible to a non-specialist reader and undetectable without checking the cited text.

Attribution research supplies the formal apparatus for that check. The Attributable to Identified Sources framework defines a generated statement as attributable only if a generic hearer would affirm ``according to source $X$, $S$'' given the cited evidence, and supplies an annotation protocol demonstrating that the property is measurable rather than intuitive~\cite{rashkin2023attribution}. Attributed question answering operationalises the criterion into a task of answer-plus-supporting-passage, and its architectural comparison finds retrieve-then-read pipelines superior to post-hoc citation attachment on both correctness and attribution~\cite{bohnet2022attributed}. That result is the justification for constraining generation to retrieved source identifiers at generation time rather than adding citations afterwards.

The scale of the problem in deployed systems is documented directly. Liu et al.\ human-audit four commercial generative search engines and find that only about half of generated sentences are fully supported by their own citations, and that roughly a quarter of citations do not support the statement they are attached to, while fluency and perceived utility correlate \emph{inversely} with citation precision~\cite{liu2023verifiability}. Citation presence and citation support are thus distinct quantities, and only the former is cheap to measure. ALCE turned this observation into a reproducible benchmark, scoring fluency, correctness, and citation precision and recall separately over retrieval corpora, and confirming that strong models leave a substantial fraction of statements uncited or wrongly cited~\cite{gao2023alce}. RARR offers the complementary repair path: rather than regenerating, retrieve evidence per claim and minimally revise unsupported spans, trading attribution against preservation of the original output~\cite{gao2023rarr}. This is the published antecedent for treating a verification failure as a trigger for targeted revision rather than as an immediate abstention.

The architecture examined in this work applies structural safeguards derived from this literature: the generator returns source identifiers rather than free-form citation strings, supporting quotes must occur verbatim in the supplied evidence, and the application resolves identifiers to canonical citations. Structural checks are deterministic and cheap, but they do not establish semantic entailment. A verbatim quote may be accurate and still fail to support the claim built on it, and material omission is invisible to a quote check. Claim support must therefore be assessed separately, and the evaluation in Chapter~\ref{chap:results} reports the two families of check independently.

\section{Selective Answering, Calibration, and Escalation}
\label{sec:sota-selective}

A system that must answer every question cannot be safe on a corpus where some questions are ambiguous and some answers are unsupported. Selective prediction provides the formal treatment: a predictor is paired with a rejection function, and performance is characterised by a risk--coverage curve rather than by a single accuracy figure, with distribution-free guarantees available at a chosen coverage~\cite{geifman2017selective}. Framing abstention this way converts an apparently ad hoc threshold into a defensible operating point.

The signal driving that decision requires care. Kamath et al.\ show that model confidence is badly miscalibrated under domain shift and that a separate calibrator trained on mixed in-domain and out-of-domain data substantially outperforms thresholding the model's own softmax~\cite{kamath2020selective}, which argues for separating the answering model from the abstention decision. Earlier work on question-answering calibration reached a compatible conclusion, finding models poorly calibrated out of the box and improvable by fine-tuning, temperature scaling, and candidate paraphrasing~\cite{jiang2021calibration}. Against this, Kadavath et al.\ show that large models are reasonably well calibrated in suitable formats, can self-evaluate sampled answers, and can be trained to predict whether they know an answer, with partial generalisation to unseen tasks~\cite{kadavath2022know}. Taken together, model-internal confidence is usable as \emph{one} input to an escalation trigger and unsafe as the only one --- which is why the policy developed here combines retrieval disagreement, score margin, and verifier outcome rather than relying on a generator score.

Abstention can also be learned directly. R-Tuning diagnoses standard instruction tuning as forcing an answer regardless of parametric knowledge, then partitions training data by what the model demonstrably knows and fine-tunes explicit refusal on the uncertain partition, improving both accuracy and abstention and showing that refusal behaves as a transferable skill~\cite{zhang2024rtuning}. This work does not train refusal, because the generator is not trainable in the deployed configuration; the result is recorded because it establishes that a rule-based abstention gate is a design choice with a learned alternative, not the only available mechanism.

\section{Evaluating Retrieval and Generation}
\label{sec:sota-eval}

Retrieval evaluation is comparatively settled. Recall at a fixed depth, \ac{MRR}, and \ac{nDCG} are standard, the last introduced to credit graded relevance and to discount gain by rank position~\cite{jarvelin2002cumulated}. BEIR established that retrieval models behave very differently across tasks and domains, so a model selected on one benchmark cannot be assumed to transfer~\cite{thakur2021beir}. KILT is the closer methodological precedent for this work: it unifies knowledge-intensive tasks over a single corpus snapshot and mandates provenance annotations, so that systems are scored on retrieved evidence as well as on final output~\cite{petroni2021kilt}.

Generation evaluation is less settled and more consequential. ROUGE measures recall-oriented n-gram overlap against reference text~\cite{lin2004rouge} and remains widely reported, but Maynez et al.\ demonstrate through large-scale human annotation that such overlap metrics correlate only weakly with human faithfulness judgements, while entailment-based measures correlate considerably better~\cite{maynez2020faithfulness}. A ROUGE score therefore cannot license a conclusion about factual correctness. This is not a marginal caveat in the present setting: the earlier phase of this project reported ROUGE-L alongside a citation-matching proxy, and the reinterpretation of those figures in Chapter~\ref{chap:results} rests directly on Maynez et al.

Purpose-built \ac{RAG} evaluation frameworks decompose quality along the pipeline. RAGAS scores faithfulness, answer relevance, and context relevance without reference answers, using prompted model judgements~\cite{es2024ragas}. ARES trains lightweight judges on synthetic data and then corrects their bias using prediction-powered inference over a small human-annotated set~\cite{saadfalcon2024ares}. ARES is the most directly applicable precedent for the design adopted here, because it addresses the same structural problem: a largely machine-generated evaluation resource made trustworthy by a human-anchored correction, with the uncertainty of that correction reported rather than suppressed.

Both frameworks depend on model-generated judgements, which imports a known set of biases. Zheng et al.\ characterise language models as judges systematically, documenting position, verbosity, and self-enhancement bias alongside weak grading of reasoning, while also showing agreement with human judgements at a level comparable to human--human agreement~\cite{zheng2023judging}. The correct inference is neither to reject automated judging nor to accept it silently, but to report the human-agreement figure that validates it in the specific setting of use.

\section{Benchmark Validity and Reviewer Agreement}
\label{sec:sota-validity}

An evaluation is only as trustworthy as its gold data, and this is where the present work makes an unusually large methodological investment. Two threats are relevant.

The first is contamination. Sainz et al.\ argue that closed pre-training corpora have absorbed the standard public benchmarks, that this invalidates measured progress, and that contamination must be documented per benchmark rather than assumed absent~\cite{sainz2023contamination}. Building a fresh benchmark over a specific regulatory instrument mitigates this threat by construction, at the cost of introducing the second one.

The second is the validity of machine-generated evaluation items. Where questions and reference answers are produced by a model from the same family that is later evaluated, agreement between system and gold may reflect shared provenance rather than correctness. The mitigation adopted in this work is human audit of a stratified sample, with agreement reported rather than asserted. That requires chance-corrected agreement statistics: Cohen's coefficient for two raters on nominal categories~\cite{cohen1960coefficient}, Fleiss's generalisation for a fixed number of raters drawn from a larger pool~\cite{fleiss1971measuring}, and the conventional verbal benchmarks for interpreting the resulting magnitudes~\cite{landis1977measurement}. Landis and Koch present those cut-offs as arbitrary conventions, and they are treated as such here.

Reported proportions require intervals. At the sample sizes typical of a human-audited subset, and at proportions near zero or one, the normal-approximation interval is inappropriate; the score-based interval obtained by inverting the normal test on the true proportion remains within the unit interval and is used throughout~\cite{wilson1927probable}.

A final observation belongs here rather than in the results. Where independent qualified reviewers disagree about whether an answer is supported by a provision, that disagreement is evidence about the question, not noise in the measurement. The literature on legal reasoning benchmarks acknowledges genuine indeterminacy in rule application~\cite{guha2023legalbench,holzenberger2020sara}, and an evaluation protocol that resolves every disagreement by majority vote discards exactly the signal that identifies ambiguous questions. This work therefore reports contested records separately.

\section{Agentic and Multi-Step Retrieval}
\label{sec:sota-agentic}

Multi-step retrieval can decompose a question, follow cross-references, or issue a revised query after inspecting initial evidence. Self-RAG demonstrates the value of learning when to retrieve at all~\cite{asai2024selfrag}, and post-hoc revision architectures show that a failed verification need not force abstention~\cite{gao2023rarr}. These capabilities are genuinely useful when a validated residual requires several sources, which is a realistic prospect in regulation, where obligations are frequently defined by cross-reference.

They also introduce latency, monetary cost, and additional failure modes, and they complicate attribution: an answer assembled over several retrieval rounds is harder to trace to a single governing provision. Applying multi-step retrieval to every question would obscure whether a simpler retriever, a reranker, a clarification request, or an abstention rule was sufficient. The position taken in this work, and recorded in the decision log, is that agentic retrieval is admitted only after a residual failure has been shown to require multi-step evidence following independent review of the benchmark. That ordering --- diagnose first, escalate second --- is the same principle the contribution advances at the level of individual questions.

\section{Positioning of This Work}
\label{sec:sota-gap}

The literature reviewed above supports each component of the pipeline individually. Hybrid retrieval is well motivated~\cite{cormack2009rrf,bruch2024fusion,chen2024m3embedding}; reranking is well established~\cite{nogueira2019reranking,nogueira2020monot5}; granularity effects are measured~\cite{chen2024densex,liu2024lostmiddle}; attribution is formalised and its failure in deployed systems is documented~\cite{rashkin2023attribution,liu2023verifiability,magesh2025hallucination}; selective answering has a statistical foundation~\cite{geifman2017selective,kamath2020selective}; and evaluation frameworks exist for each stage~\cite{petroni2021kilt,es2024ragas,saadfalcon2024ares}.

Three gaps remain, and together they define this work.

First, these components are usually evaluated in isolation and on general-domain corpora. There is no published study that assembles them over a single European aviation-regulation instrument and reports, for the same questions, where each stage succeeds and fails.

Second, the interventions are typically presented as improvements to be applied uniformly. The question of \emph{when} an intervention is warranted --- which observable runtime condition justifies reranking, expansion, clarification, or human referral for a particular question --- is largely unaddressed. Self-RAG learns this control internally~\cite{asai2024selfrag}; a compliance setting needs it externalised and inspectable.

Third, and most consequentially, the field reports system failure against gold data whose own validity is rarely established. Where an apparent retrieval ceiling is in fact a mixture of genuine system failure, incomplete gold evidence, and irreducibly ambiguous questions, reporting it as a ceiling is a measurement error rather than a finding.

\begin{table}[htbp]
  \centering
  \caption{Escalation levels and the prior work that motivates each. The contribution of this thesis is the calibrated transition between levels, not the levels themselves.}
  \label{tab:escalation-literature}
  \small
  \begin{tabularx}{\textwidth}{@{}c l X@{}}
    \toprule
    \textbf{Level} & \textbf{Intervention} & \textbf{Motivating prior work} \\
    \midrule
    0 & Dense retrieval over flat chunks & \cite{lewis2020rag,karpukhin2020dpr} \\
    1 & Hierarchy-aware parent--child retrieval & \cite{chen2024densex,liu2024lostmiddle} \\
    2 & Hybrid lexical and dense candidates with fusion & \cite{robertson2009bm25,cormack2009rrf,bruch2024fusion} \\
    3 & Cross-encoder reranking of the candidate pool & \cite{nogueira2019reranking,nogueira2020monot5,wang2020minilm} \\
    4 & Evidence-constrained generation and verification & \cite{rashkin2023attribution,bohnet2022attributed,gao2023alce} \\
    5 & Clarification or targeted revision & \cite{gao2023rarr} \\
    6 & Abstention and referral to human review & \cite{geifman2017selective,kamath2020selective,magesh2025hallucination} \\
    \bottomrule
  \end{tabularx}
\end{table}

The contribution of this thesis is accordingly not a new foundation model, nor a new retrieval architecture. It is a controlled, evidence-first study of where a regulatory \ac{RAG} pipeline fails and how a system should respond, conducted under a protocol that establishes the validity of its own gold data before quoting any ceiling. The work combines hierarchy-aware corpus construction, hybrid retrieval, local reranking, explicit evidence contracts, audited benchmark construction with reported reviewer agreement, and selective escalation in one reproducible protocol. Its evaluation separates data defects from retrieval, ranking, generation, and verification failures, so that each intervention is justified by a diagnosed failure mode rather than adopted because it is available.
